\begin{abstract}
We report the observation of accelerator-produced neutrino interactions in the ProtoDUNE Horizontal Drift liquid argon time projection chamber at CERN, operated parasitically during 2024 North Area running. Located downstream of the SPS 400 GeV proton beam and T2 target, ProtoDUNE functions as an opportunistic beam-dump experiment sensitive to weakly interacting particles. Using a dedicated trigger and offline selection, we identify an excess of events aligned with the SPS/T2 beam direction, in accord with simulation predictions.
These results establish a proof of principle for beam-induced neutrino detection at ProtoDUNE and motivate future new-physics searches.
\end{abstract}